Die mittlere Schnelligkeit der Teilchen eines idealen Gases ist $\bar v = \sqrt{8kT/(\pi m)}$ (FoTaBe S.~171). Mit der Teilchenmasse $m = M/N_\text{A}$ und $R = N_\text{A}k$ wird daraus $\bar v = \sqrt{8RT/(\pi M)}$.

Die molare Masse von Butan ist (FoTaBe S.~228)
\begin{align}
 M &= \MBuF = 4\cdot\MCaO + 10\cdot\MHaO = \MBug
\end{align}
Mit $T = \TBuF = \TBu$ wird
\begin{align}
 \bar v &= \vBuF \\
   &= \sqrt{\frac{8\cdot\ncR\cdot\TBu}{\pi\cdot\MBu}} \\
   &= \vBu \approx \result{\vBuP{0}{2}}
\end{align}
Die Wurzel aus dem mittleren Geschwindigkeitsquadrat ist etwas grösser:
\begin{align}
 v_\text{rms} &= \vrBuF = \vrBu \approx \vrBuP{0}{2}
\end{align}
